\subsection{Fixed-schedule comparison with explicit high-probability rates}
\label{subsec:fixed-schedule-package}

We assemble the deterministic ODE--jump comparison and the two
stochastic concentration bounds into the final fixed-schedule tracking theorem.

This subsection packages the detailed estimates from
\cref{lemma-common-empirical-moment-residual-directional-rt},
\cref{lem:jump-bound-common-seed},
\cref{lem:piecewise-gronwall-forced-jumps},
\cref{lemma-uniform-sup-gamma-bound-convenient},
\cref{cor:tilde-gamma-jump-time-union-bound}
into a single fixed-schedule statement.  

We keep only the quantities that enter the final comparison bound.
\begin{itemize}
    \item the pathwise comparison variable is
    \[
    r_t=\frac{1}{\sqrt N}\norm{\bs X_t-\bs Y_t}_2,
    \]
    \item the continuous stochastic residual is
    \[
    \Gamma_t^N=\sum_{m\ge1}\mathrm L_Da_m
    \left\|\frac1N\sum_{i=1}^N\phi_m(Y_t^i)
                         -\int\phi_m(y)\,\mu_t(dy)\right\|_{\mathsf H_m},
    \]
    as in \eqref{eq:common-moment-residual-and-lipschitz-sum},
    \item the jump defect is controlled by $\widetilde\Gamma_{t_\ell-}^N$
    from \eqref{eq:tilde-gamma-full-definition-weighted} only at the prescribed jump times,
    \item curvature is retained through the propagator
    $G_\kappa^Y(t,s)=\exp(-\int_s^t\kappa_u^{Y,N}\,du)$, using the
    sampled full Hessian curvature in \eqref{eq:curvature-main}.
\end{itemize}

The first lemma solves the deterministic comparison inequality across ODE
intervals and multiplicative jump errors.

\begin{lemma}[Piecewise Gr\"onwall estimate with recursive jump amplification]
\label{lem:piecewise-gronwall-forced-jumps}
Fix $T>0$ and jump times
\[
0<t_1<\cdots<t_{\NofJ}\le T,
\qquad
t_0:=0,
\qquad
t_{\NofJ+1}:=T.
\]
Assume that $r_t:=\frac{1}{\sqrt N}\norm{\bs X_t-\bs Y_t}_2$
is absolutely continuous on each interval $(t_\ell,t_{\ell+1})$, $\ell=0,\dots,\NofJ$.
Assume right-continuity and absolute continuity up to each branch endpoint. Here $c_D$ is a deterministic upper bound on the ODE quadratic coefficient, and $q(p)=1+\sqrt{2\log(2/p)}$. For a given pathwise radius $R>0$, define
\[
\tau_R:=\inf\{t\in[0,T]: r_t\ge R\},
\]
with the convention $\inf\emptyset=\infty$.

Let
\[
\beta_t:=\kappa_t^{Y,N}-c_D R-\Gamma_t^N.
\]
Assume the following.

\begin{enumerate}
    \item For every $\ell=0,\dots,\NofJ$ and for a.e.
    $t\in (t_\ell,t_{\ell+1})\cap[0,\tau_R)$,
    \begin{equation}
    \label{eq:forced-ode-assumption-recursive}
    \dot r_t\le -\beta_t r_t+\Gamma_t^N.
    \end{equation}
    (corresponding to the conclusion in \cref{lemma-common-empirical-moment-residual-directional-rt})

    \item For each jump time $t_\ell\le\tau_R$, including a candidate first exit at that jump, $\ell=1,\dots,\NofJ$, the previous jump lemma yields
    \begin{equation}
    \label{eq:jump-recursion-assumption}
    r_{t_\ell}
    \le
    \bigl(1+a_\ell\bigr)r_{t_\ell-}+b_\ell,
    \end{equation}
    where
    \[
    a_\ell(\supXY,p):=c_J(D)\eta_\ell q(p),
    \qquad
    b_\ell:=a_\ell\Bigl(\sqrt{\Gamma_{t_\ell-}^N}+\sqrt{\widetilde{\Gamma}_{t_\ell-}^N}\Bigr).
    \]
    (corresponding to \cref{lem:jump-bound-common-seed} on the event $E_{T,\supXY}$).
\end{enumerate}

Then, for every $t\in[0,T]\cap[0,\tau_R)$,
\begin{align}
\label{eq:piecewise-recursive-bound}
r_t
&\le
\int_0^t
\exp\Bigl(-\int_s^t \beta_u\,du\Bigr)
\prod_{\ell:\,s<t_\ell\le t}(1+a_\ell)\,
\Gamma_s^N\,ds\\
&+
\sum_{\ell:\,t_\ell\le t}
\exp\Bigl(-\int_{t_\ell}^t \beta_u\,du\Bigr)
\prod_{j:\,t_\ell<t_j\le t}(1+a_j)\,
b_\ell
\end{align}
if $r_0=0$. The same bound holds at a finite candidate exit $\tau_R$, using the pre-jump or post-jump trace as appropriate. No adaptedness of $R$ is required.
\end{lemma}

\begin{proof}
\noindent\textbf{Step 1: integrate within one ODE interval.}
For $0\le s\le t\le T$, define
\[
G(t,s):=\exp\Bigl(-\int_s^t \beta_u\,du\Bigr).
\]
By the definition of $\beta$, this is
$G(t,s)=G_\kappa^Y(t,s)\exp\{c_DR(t-s)+\int_s^t\Gamma_u^N\,du\}$,
with the curvature propagator of \cref{def:curvature-responses}.
On each continuity interval $(t_\ell,t_{\ell+1})$, the differential inequality
\eqref{eq:forced-ode-assumption-recursive} implies
\begin{equation}
\label{eq:continuous-duhamel-recursive}
r_t
\le
G(t,s)\,r_s+\int_s^t G(t,u)\Gamma_u^N\,du,
\qquad
t_\ell\le s\le t<t_{\ell+1},
\quad t<\tau_R.
\end{equation}
Indeed,
\[
\dot r_t+\beta_t r_t\le \Gamma_t^N,
\]
and multiplying by the integrating factor
$\exp(\int_s^t\beta_u\,du)$ and integrating from $s$ to $t$
gives \eqref{eq:continuous-duhamel-recursive}.


\par\noindent\textbf{Step 2: propagate the bound through each jump and interval.}
We now prove \eqref{eq:piecewise-recursive-bound} by induction over the jump intervals.

First, if $t\in[0,t_1)\cap[0,\tau_R)$, then there are no jumps before time $t$, hence
\eqref{eq:piecewise-recursive-bound} reduces exactly to
\eqref{eq:continuous-duhamel-recursive} with $s=0$.

Next, let $m\in\{1,\dots,\NofJ\}$ and assume that
\eqref{eq:piecewise-recursive-bound} has been proved for all times in $[0,t_m)$.
Take $t\in[t_m,t_{m+1})\cap[0,\tau_R)$.
Applying \eqref{eq:continuous-duhamel-recursive} on $[t_m,t]$ yields
\begin{equation}
\label{eq:after-last-jump-recursive}
r_t
\le
G(t,t_m)\,r_{t_m}
+
\int_{t_m}^t G(t,s)\Gamma_s^N\,ds.
\end{equation}
Using the jump recursion \eqref{eq:jump-recursion-assumption},
\[
r_{t_m}\le (1+a_m)r_{t_m-}+b_m,
\]
we obtain
\begin{equation}
\label{eq:after-substitute-jump-recursive}
r_t
\le
G(t,t_m)(1+a_m)r_{t_m-}
+
G(t,t_m)b_m
+
\int_{t_m}^t G(t,s)\Gamma_s^N\,ds.
\end{equation}

Now apply the induction hypothesis at time $t_m-$:
\begin{align}
r_{t_m-}\le&
G(t_m,0)\prod_{\ell=1}^{m-1}(1+a_\ell)\,r_0\quad
+
\int_0^{t_m}
G(t_m,s)\prod_{\ell:\,s<t_\ell<t_m}(1+a_\ell)\Gamma_s^N\,ds\\
&+
\sum_{\ell=1}^{m-1}
G(t_m,t_\ell)\prod_{j:\,t_\ell<t_j<t_m}(1+a_j)\,b_\ell.
\label{eq:IH-at-tmminus}
\end{align}
Substituting \eqref{eq:IH-at-tmminus} into
\eqref{eq:after-substitute-jump-recursive}, we get
\begin{align*}
r_t
&\le
G(t,t_m)(1+a_m)G(t_m,0)\prod_{\ell=1}^{m-1}(1+a_\ell)\,r_0
\\
&\quad
+
G(t,t_m)(1+a_m)
\int_0^{t_m}
G(t_m,s)\prod_{\ell:\,s<t_\ell<t_m}(1+a_\ell)\Gamma_s^N\,ds
\\
&\quad
+
G(t,t_m)(1+a_m)
\sum_{\ell=1}^{m-1}
G(t_m,t_\ell)\prod_{j:\,t_\ell<t_j<t_m}(1+a_j)\,b_\ell
\\
&\quad
+
G(t,t_m)b_m
+
\int_{t_m}^t G(t,s)\Gamma_s^N\,ds.
\end{align*}
Using the semigroup property of $G$,
\[
G(t,t_m)G(t_m,s)=G(t,s),
\qquad 0\le s\le t_m\le t,
\]
we rewrite this as
\begin{align*}
r_t
&\le
G(t,0)\prod_{\ell=1}^{m}(1+a_\ell)\,r_0
+
\int_0^{t_m}
G(t,s)\prod_{\ell:\,s<t_\ell\le t_m}(1+a_\ell)\Gamma_s^N\,ds
\\
&\quad
+
\sum_{\ell=1}^{m-1}
G(t,t_\ell)\prod_{j:\,t_\ell<t_j\le t_m}(1+a_j)\,b_\ell
+
G(t,t_m)b_m
+
\int_{t_m}^t G(t,s)\Gamma_s^N\,ds.
\end{align*}
Since there is no jump in $(t_m,t)$, for $s\in(t_m,t)$ the product
$\prod_{\ell:\,s<t_\ell\le t}(1+a_\ell)$ is empty and equals $1$.
Also, because $t\in[t_m,t_{m+1})$, the sets
\[
\{\ell:\,s<t_\ell\le t\}
\quad\text{and}\quad
\{\ell:\,s<t_\ell\le t_m\}
\]
coincide when $s<t_m$, and similarly
\[
\{j:\,t_\ell<t_j\le t\}
=
\{j:\,t_\ell<t_j\le t_m\}
\qquad (\ell\le m-1),
\]
while for $\ell=m$ the product over $j:\,t_m<t_j\le t$ is empty and equals $1$.
Therefore,
\begin{align*}
r_t&\le
G(t,0)\prod_{\ell:\,t_\ell\le t}(1+a_\ell)\,r_0
+
\int_0^t
G(t,s)\prod_{\ell:\,s<t_\ell\le t}(1+a_\ell)\Gamma_s^N\,ds\\
&+
\sum_{\ell:\,t_\ell\le t}
G(t,t_\ell)\prod_{j:\,t_\ell<t_j\le t}(1+a_j)\,b_\ell,
\end{align*}
which is exactly \eqref{eq:piecewise-recursive-bound}. This completes the induction. Taking the left limit gives the bound at a continuous candidate exit; at a candidate exit that is a jump, apply the assumed recursion at that jump itself. Empty sums and products cover the case of no jumps.
\end{proof}

The next lemma packages the local coefficients without discarding their dependence on the population moment scale or anti-spikiness.

\begin{lemma}[Explicit local comparison coefficients]
\label{lem:comparison-coefficients-v15}
Fix the deterministic population scale $K=K_{\rm pop}$ and radius $D\ge1$.
Use $\mathrm L_D,\gamma_m(D)=\mathrm L_Da_m$ from
\cref{lemma-localized-moment-smoothness-consequences} and put
\[
\Lambda(D)=\sum_{m\ge1}\gamma_m(D)L_m(1+2D^{m-1}),\qquad
c_D=\Lambda(D)+(2+A_{\rm sp})\mathrm L_D,
\]
\[
c_J(D)=1+\sqrt d\{\Lambda(D)+2\mathrm L_D\}+\sqrt{2d\mathrm L_D}.
\]
On $E_{\rm pop}\cap E_{T,D}\cap E_{T,\rm sp}(A_{\rm sp})$,
$c_D$ and $c_J(D)$ are the coefficients in
\cref{lemma-common-empirical-moment-residual-directional-rt,lem:jump-bound-common-seed}.
With $p_0=p/8$, write $\Xi=8N(e+T)(\NofJ+1)(1+B)/p$ and choose
\begin{equation}
\label{eq:def-supXY-v15}
\supXY=D=C_D(1+\log\Xi)^{1/\alpha},
\end{equation}
where the fixed $C_D$ is sufficiently large for
\eqref{eq:gamma-radius-choice} with budget $p_0$ and for the clipping-bias
estimate in Lemma~\ref{lemma-uniform-sup-gamma-bound-convenient}.
For fixed model parameters and uniformly bounded $K_0,B,A_{\rm sp}$,
for every $\varepsilon>0$ there is $C_\varepsilon$ such that
\begin{equation}
\label{eq:Ccmp-localized-v15}
1+c_D+c_J(D)+\Lambda(D)\le C_\varepsilon\Xi^\varepsilon.
\end{equation}
The constants may depend on the fixed anchor derivatives and $C_3$.
\end{lemma}
\begin{proof}
The empirical moment displacement is
$L_m(1+2D^{m-1})r_t$ by \cref{lemma-directional-measure-slot-error},
\eqref{eq:paired-empirical-moment-displacement}. The sum of the two Hessian measure errors is bounded
with weights $\gamma_m(D)$ by \cref{eq:local-two-hessians-moment-bound},
so its coefficient is $\Lambda(D)$.
The two spatial kernel slots contribute $2\mathrm L_D$ by
\cref{eq:local-kernel-spatial-bound}; the diagonal term contributes
$A_{\rm sp}\mathrm L_D$ by \cref{eq:local-jacobian-spatial-bound} and
\cref{eq-cubic-antispike-bound}. Thus
the ODE coefficient is exactly $c_D$. The jump
coefficient is exactly the one in Lemma~\ref{lem:jump-bound-common-seed}.
For the common growth bound, $D\ge1$ and
\eqref{eq:conditional-comparison-envelope} give
\[
\begin{aligned}
 \Lambda(D)
 &\le2\mathrm L_D\sum_m a_mL_m
       (2+D^m+K_{\rm pop}^m m^{m/\alpha})
 \le\mathrm L_D^2.
\end{aligned}
\]
For the last inequality, use $2u\le4(1+u)^3$ for $u\ge0$ and
$1+g_0+h_0+C_3\ge1$ in the definition of $\mathrm L_D$.
Using $\mathrm L_D\ge1$ in the displayed definitions therefore yields
\begin{equation}
\label{eq:comparison-coefficients-common-envelope}
 1+c_D+c_J(D)+\Lambda(D)
 \le(6+A_{\rm sp}+3\sqrt d+\sqrt{2d})\mathrm L_D^2.
\end{equation}
By \cref{lem:population-moment-event},
\cref{eq:population-exponential-amplitude,eq:population-exponential-scale},
the grid count is at most $C(e+T)(\NofJ+1)(1+B)$, so
$\log A_0\le C+\log\Xi$ and
$K_{\rm pop}\le C(1+\log\Xi)^{1/\alpha}$ when $K_0,B$ are bounded.
Insert this and \eqref{eq:def-supXY-v15} into
\cref{lem:conditional-local-coefficient-growth}, \eqref{eq:local-envelope-growth}:
\[
 \mathrm L_D^2\le C\exp\{C[1+\log\log(e+\Xi)]^{\kappa_{\rm dec}/(\kappa_{\rm dec}-1)}\}
 \le C_\varepsilon\Xi^\varepsilon.
\]
The last step uses $[\log\log(e+\Xi)]^{\kappa_{\rm dec}/(\kappa_{\rm dec}-1)}
=o(\log\Xi)$. This proves \eqref{eq:Ccmp-localized-v15}.
\end{proof}

\paragraph{Setup and curvature-response factors for the tracking theorem.}
Fix $N\ge2$, $p\in(0,1)$, a deterministic schedule
$0<t_1<\cdots<t_{\NofJ}\le T$ with $\eta_\ell\ge0$, and $A_{\rm sp}\ge1$.
Set $p_0=p/8$ and use the population constants of
\cref{lem:population-moment-event} with budget $p_0$.
Choose $D$ satisfying \eqref{eq:gamma-radius-choice} with budget $p_0$;
\eqref{eq:def-supXY-v15} is one such choice. Use $c_D,c_J$ from
\cref{lem:comparison-coefficients-v15}.
Define the actual deterministic residual envelopes by
\begin{equation}
\label{eq:def-delta-v15}
\delta_{N,T,p}:=\delta_\Gamma(D)
\quad\text{from \eqref{eq:gamma-explicit-envelope}, with }p_\Gamma=p_0,
\end{equation}
\begin{equation}
\label{eq:def-tdelta-v15}
\widetilde\delta_{N,T,p}:=
\widetilde\delta_{N,T,p_0}^{\rm HS}
\quad\text{from \eqref{eq:phase1-HS-tilde-delta-envelope}},
\end{equation}
For readability write $\delta=\delta_{N,T,p}$ and
$\widetilde\delta=\widetilde\delta_{N,T,p}$ in the following statement and proof.

For the sampled full Hessian curvature $\kappa_t^{Y,N}$, on any path
where $\int_0^T|\kappa_t^{Y,N}|\,dt<\infty$, recall the curvature-response factors
of \cref{def:curvature-responses}:
\begin{equation}
\label{eq:tracking-curvature-responses}
G_\kappa^Y(t,s)=\exp\!\left(-\int_s^t\kappa_u^{Y,N}\,du\right),\qquad
\mathfrak C_t^Y=\int_0^tG_\kappa^Y(t,s)\,ds,\qquad
\mathfrak J_t^Y=\sum_{\ell:t_\ell\le t}\eta_\ell G_\kappa^Y(t,t_\ell).
\end{equation}
Set
\[
q_p=1+\sqrt{2\log(2(\NofJ+1)/p_0)},\qquad
a_\ell=c_J(D)q_p\eta_\ell,
\]
\begin{equation}
\label{eq:def-jump-amplification-v15}
J_{N,T,p}:=\prod_{\ell=1}^{\NofJ}(1+a_\ell),
\end{equation}
\begin{equation}
\label{eq:def-B-v15}
\mathfrak B_{N,T,p}(t):=
\delta\mathfrak C_t^Y+
c_J(D)q_p(\sqrt\delta+\sqrt{\widetilde\delta})\mathfrak J_t^Y,
\qquad \mathfrak B_*:=\sup_{t\le T}\mathfrak B_{N,T,p}(t).
\end{equation}
Empty sums and products are zero and one, respectively. The envelope
$\mathfrak B$ is generally random because the curvature-response factors are random.
Let $R=2J_{N,T,p}\mathfrak B_*$ and define the nonlinear stability event
\begin{equation}
\label{eq:smallness-v15}
\mathcal S_{N,T,p}:=
\{T(\delta+c_D R)\le1/4\}.
\end{equation}
\paragraph{Which quantities are fixed.}
The following convention applies to the asymptotic corollary; the explicit
theorem retains the actual deterministic envelopes.
\begin{center}
\begin{tabular}{p{0.25\linewidth}p{0.66\linewidth}}
Fixed structural bounds & $d,\alpha$, feature growth/decay parameters,
$g_0,h_0,C_3$, and the initial population law.\\
Uniform regularity assumptions & $K_0,B,A_{\rm sp}$; these are assumptions
along the growing horizon, not consequences of localization.\\
Deterministic growing quantities & $D,K_{\rm pop},\delta_{N,T,p},
\widetilde\delta_{N,T,p},c_D,c_J,\Lambda,q_p,J_{N,T,p}$.
The residual envelopes are defined in
\cref{eq:def-delta-v15,eq:def-tdelta-v15}; their dependence on $N,T,p$ is retained.\\
Random reference quantities & $\kappa^{Y,N},\mathfrak C^Y,\mathfrak J^Y,
\mathfrak B_*,R$ and the nonlinear stability event. Their bounds are pathwise conditions.
\end{tabular}
\end{center}

\begin{theorem}[Fixed-schedule comparison on the nonlinear stability event]
\label{thm:fixed-schedule-high-prob-r-v15}
Assume \cref{assumption-structured-moment-smoothness,assumption-common-seed,assumption-moment-growth,assumption-uniform-concentration-piecewise-lipschitz}
and use the deterministic schedule, $N\ge2$, $p\in(0,1)$, and radius
specified above, with $r_0=0$ and integrable curvature on $[0,T]$.
Use the residual envelopes \cref{eq:def-delta-v15,eq:def-tdelta-v15},
the amplification \eqref{eq:def-jump-amplification-v15}, the forcing
\eqref{eq:def-B-v15}, and the nonlinear stability event \eqref{eq:smallness-v15}.
There is an event $\Omega_p$ with $\mathbb P(\Omega_p)\ge1-p$ such that on
$\Omega_p\cap E_{T,\rm sp}(A_{\rm sp})\cap\mathcal S_{N,T,p}$,
simultaneously at all times and left limits,
\begin{equation}
\label{eq:high-prob-r-bound-v15}
r_t\le e^{1/4}J_{N,T,p}\mathfrak B_{N,T,p}(t)
<\tfrac32J_{N,T,p}\mathfrak B_{N,T,p}(t)
\quad\text{when }\mathfrak B_{N,T,p}(t)>0,
\end{equation}
and $r_t=0$ if that envelope is zero. In particular,
\[
\mathbb P\!\left(E_{T,\rm sp}(A_{\rm sp})\cap\mathcal S_{N,T,p}
\cap\{\exists t\le T:r_t>2J_{N,T,p}\mathfrak B_{N,T,p}(t)\}\right)\le p.
\]
This is a bound on failure within the stated events, not a bound on the
conditional failure probability given $\mathcal S_{N,T,p}$.
\end{theorem}
\par\noindent\textbf{Remark (Relation to the main-text growing-horizon theorem).}
The main-text result in \Cref{theorem-main-r_t} applies the preceding theorem
on paths satisfying the controlled-curvature path condition $\mathcal R_N$ and uses the simple sufficient horizon
$T\le N^{1/4-\upsilon-\zeta}$. The integrable-curvature condition imposes no separate
model assumption there: on the common good event used in the proof,
\Cref{eq:fixed-schedule-curvature-bound} gives
$|\kappa_t^{Y,N}|\le2\mathrm L_D$, and hence
\[
 \int_0^T|\kappa_t^{Y,N}|\,\dd t
 \le \int_0^T2\mathrm L_D\,\dd t
 =2T\mathrm L_D<\infty.
\]
The explicit stability requirement remains
$T(\delta_{N,T,p}+c_DR)\le1/4$ by
\Cref{eq:smallness-v15}; it can also hold on longer stable trajectories. The
strict margin $\zeta>0$ in the main-text sufficient regime absorbs the
subpolynomial concentration and localization factors. Finally, the probability
statement is a bound on failure within the controlled-curvature path condition
and the nonlinear stability event. It neither lower-bounds
$\mathbb P(\mathcal R_N)$ nor gives a conditional failure probability given
that condition.
\begin{proof}
\noindent\textbf{Step 1: a common good event.}
Recall $p_0=p/8$ and the residual envelopes
\eqref{eq:def-delta-v15}--\eqref{eq:def-tdelta-v15}. Define
\[
 E_\Gamma:=\{\sup_{t\le T}\Gamma_t^N\le\delta\},\qquad
 E_{\widetilde\Gamma}:=
 \{\max_{\ell\le\NofJ}\widetilde\Gamma_{t_\ell-}^N\le\widetilde\delta\}.
\]
Every supremum includes the relevant left and right traces. Use the
population and radius events with budgets $p_0$ specified in the statement.
For each jump, let $\Omega_{\ell,p_0/(\NofJ+1)}^{\rm jump}$ be the Gaussian
event in \cref{lem:jump-bound-common-seed}, and set
\[
 \Omega_p:=E_{\rm pop}\cap E_{T,D}\cap E_\Gamma\cap E_{\widetilde\Gamma}
       \cap\bigcap_{\ell=1}^{\NofJ}
                   \Omega_{\ell,p_0/(\NofJ+1)}^{\rm jump}.
\]
The failure bounds and their sources are:
\begin{center}
\begin{tabular}{p{0.38\linewidth}p{0.10\linewidth}p{0.42\linewidth}}
Failure event & Budget & Source\\\hline
$E_{\rm pop}^c$ & $p_0$ & \cref{lem:population-moment-event}\\
$E_{T,D}^c$ & $p_0$ & \cref{lemma:uniform-radius-bound-XY}, \eqref{eq:gamma-radius-choice}\\
$E_{\rm pop}\cap E_{T,D}\cap E_\Gamma^c$ & $p_0$ & \cref{lemma-uniform-sup-gamma-bound-localized-conditional}, \eqref{eq:gamma-localized-probability}\\
$E_{\rm pop}\cap E_{\widetilde\Gamma}^c$ & $p_0$ & \cref{cor:tilde-gamma-jump-time-union-bound}, \eqref{eq:tilde-gamma-jump-time-union-bound}\\
$(\Omega_{\ell,p_0/(\NofJ+1)}^{\rm jump})^c$ & $\frac{p_0}{\NofJ+1}$ & \cref{lem:jump-bound-common-seed}
\end{tabular}
\end{center}
The covariance estimate conditions on $\mathcal G_{\rm pop}$, while the
Gaussian jump estimate conditions on $\mathcal H_{\ell-}$. The tower
property converts these bounds to unconditional ones:
\[
\begin{aligned}
 \mathbb P(E_{\rm pop}\cap E_{\widetilde\Gamma}^c)
 &=\mathbb E[\mathbf1_{E_{\rm pop}}
      \mathbb P(E_{\widetilde\Gamma}^c\mid\mathcal G_{\rm pop})]\le p_0,\\
 \mathbb P((\Omega_{\ell,p_0/(\NofJ+1)}^{\rm jump})^c)
 &=\mathbb E[\mathbb P((\Omega_{\ell,p_0/(\NofJ+1)}^{\rm jump})^c
      \mid\mathcal H_{\ell-})]\le p_0/(\NofJ+1).
\end{aligned}
\]
Consequently,
\[
\begin{aligned}
 \mathbb P(\Omega_p^c)
 &\le\mathbb P(E_{\rm pop}^c)+\mathbb P(E_{T,D}^c)
      +\mathbb P(E_{\rm pop}\cap E_{T,D}\cap E_\Gamma^c)\\
 &\quad+\mathbb P(E_{\rm pop}\cap E_{\widetilde\Gamma}^c)
      +\sum_{\ell=1}^{\NofJ}
              \mathbb P((\Omega_{\ell,p_0/(\NofJ+1)}^{\rm jump})^c)\\
 &\le4p_0+\frac{\NofJ}{\NofJ+1}p_0<5p_0<p.
\end{aligned}
\]
This union bound does not require independence. The sample event $E_{T,D}$
is intersected after concentration; it is not part of the conditioning used
to establish the reference product law.

\par\noindent\textbf{Step 2: deterministic comparison on the good event.}
Fix a realization in $\Omega_p\cap E_{T,\rm sp}(A_{\rm sp})$.
The two deterministic inputs can now be tied to explicit parts of this event.
On
$E_{\rm pop}\cap E_{T,D}\cap E_{T,\rm sp}(A_{\rm sp})\cap E_\Gamma$,
\cref{lemma-common-empirical-moment-residual-directional-rt} gives
\[
 \dot r_t\le-\kappa_t^{Y,N}r_t+c_Dr_t^2+\delta r_t+\delta.
\]
On the other hand, on
$E_{\rm pop}\cap E_{T,D}\cap E_\Gamma\cap E_{\widetilde\Gamma}
\cap\Omega_{\ell,p_0/(\NofJ+1)}^{\rm jump}$,
\cref{lem:jump-bound-common-seed} gives
\[
 r_{t_\ell}\le(1+a_\ell)r_{t_\ell-}
                 +a_\ell(\sqrt\delta+\sqrt{\widetilde\delta}).
\]
Before applying these two inputs,
\cref{lemma-localized-moment-smoothness-consequences},
\eqref{eq:local-field-magnitude-bound}, and the block estimate in
\cref{lemma-operator-norm-block-matrix} give
\begin{equation}
\label{eq:fixed-schedule-curvature-bound}
 |\kappa_t^{Y,N}|\le\|\HessianMatrices_t^Y\|_{\mathbb R^{dN}\to\mathbb R^{dN}}
 \le\left\|(\mathrm L_D/N)\mathbf1\mathbf1^\top+\mathrm L_DI_N
        \right\|_{\mathbb R^N\to\mathbb R^N}=2\mathrm L_D.
\end{equation}
Here $\mathbf1=(1,\ldots,1)^\top\in\mathbb R^N$ and
$P=N^{-1}\mathbf1\mathbf1^\top$ satisfies $P^\top=P$, $P^2=P$,
and $P\mathbf1=\mathbf1$. Thus $\|P+I_N\|_{\mathbb R^N\to\mathbb R^N}=2$.
In particular $\int_0^T|\kappa_t^{Y,N}|\,dt\le2T\mathrm L_D<\infty$
on this good event, so the curvature-response factors in
\eqref{eq:tracking-curvature-responses} are well-defined.
The absolute-continuity premise of
\cref{lem:piecewise-gronwall-forced-jumps} is also automatic on each ODE
branch. Indeed, the branchwise $B$-Lipschitz bounds in
\cref{assumption-uniform-concentration-piecewise-lipschitz} give
\[
\begin{aligned}
 |r_t-r_s|
 &\le\frac1{\sqrt N}\|\Delta_t-\Delta_s\|_{\mathbb R^{dN}}\\
 &\le\frac1{\sqrt N}
       \left[\sum_{i=1}^N
       \bigl(\|X_t^i-X_s^i\|_{\mathbb R^d}
             +\|Y_t^i-Y_s^i\|_{\mathbb R^d}\bigr)^2\right]^{1/2}
 \le2B|t-s|.
\end{aligned}
\]
Suppose first $\mathfrak B_*>0$, and recall
$R=2J_{N,T,p}\mathfrak B_*$ from the statement. Before a first exit from
$r<R$, replace $c_Dr_t^2$ by $c_DRr_t$.
The exponent in \cref{lem:piecewise-gronwall-forced-jumps} satisfies
\[
\begin{aligned}
 \exp\left(-\int_s^t\beta_u\,du\right)
 &=G_\kappa^Y(t,s)
       \exp\left((t-s)c_DR+\int_s^t\Gamma_u^N\,du\right)\\
 &\le G_\kappa^Y(t,s)e^{T(c_DR+\delta)}.
\end{aligned}
\]
Every indicated subproduct of $(1+a_\ell)$ is at most
$J_{N,T,p}=\prod_{\ell=1}^{\NofJ}(1+a_\ell)$, since $a_\ell\ge0$.
Apply \cref{lem:piecewise-gronwall-forced-jumps},
\eqref{eq:piecewise-recursive-bound}, retaining the two forcing terms:
\begin{equation}
\label{eq:tracking-stopped-propagator-bound}
 r_t\le J_{N,T,p}e^{T(c_DR+\delta)}
 \left[\int_0^tG_\kappa^Y(t,s)\Gamma_s^N\,ds
       +\sum_{\ell:t_\ell\le t}G_\kappa^Y(t,t_\ell)b_\ell\right].
\end{equation}
The continuous and jump forcing terms satisfy
\[
\begin{aligned}
 \int_0^tG_\kappa^Y(t,s)\Gamma_s^N\,ds
 &\le\delta\mathfrak C_t^Y,\\
 \sum_{\ell:t_\ell\le t}G_\kappa^Y(t,t_\ell)b_\ell
 &\le c_J(D)q_p(\sqrt\delta+\sqrt{\widetilde\delta})
                \sum_{\ell:t_\ell\le t}\eta_\ell G_\kappa^Y(t,t_\ell)\\
 &=c_J(D)q_p(\sqrt\delta+\sqrt{\widetilde\delta})\mathfrak J_t^Y.
\end{aligned}
\]
Here $b_\ell$ is the additive forcing in
\eqref{eq:jump-recursion-assumption}. The sum of the two right-hand sides is
$\mathfrak B_{N,T,p}(t)$ from \eqref{eq:def-B-v15}. Substitution into
\eqref{eq:tracking-stopped-propagator-bound} gives
\begin{equation}
\label{eq:tracking-forcing-envelope-bound}
 r_t\le J_{N,T,p}e^{T(c_DR+\delta)}\mathfrak B_{N,T,p}(t).
\end{equation}

\par\noindent\textbf{Step 3: exclude all first exits.}
The candidate-exit conclusion of
\cref{lem:piecewise-gronwall-forced-jumps},
\eqref{eq:piecewise-recursive-bound}, applies at a continuous exit,
a pre-jump trace, or a post-jump exit, including a jump at $T$.
At a post-jump exit its partial products already include that jump.
Thus the same bounds on the propagator, products, and forcing give
\cref{eq:tracking-stopped-propagator-bound,eq:tracking-forcing-envelope-bound}
at the candidate exit itself; no extra factor $1+a_\ell$ is introduced.
On the nonlinear stability event \eqref{eq:smallness-v15}, the candidate value is
therefore at most
\[
 e^{1/4}J_{N,T,p}\mathfrak B_*
 =\frac{e^{1/4}}2 R<R.
\]
There are finitely many branches and jumps, so this strict margin excludes
an exit on $[0,T]$. The radius may depend on future reference curvature-response factors:
the argument is pathwise and uses no stopping-time theorem.

If $\mathfrak B_*=0$ and $T>0$, positivity of $G_\kappa^Y(t,s)$ implies
$\delta=0$, and every nonnegative jump forcing
$a_\ell(\sqrt\delta+\sqrt{\widetilde\delta})$ is zero. On a branch and
below any fixed positive radius $R'$, the ODE inequality becomes
$\dot r_t\le(|\kappa_t^{Y,N}|+c_DR')r_t$.
Multiplication by the integrating factor, starting from $r=0$, gives
$r_t\le0$ on that branch. The jump recursion sends zero to zero.
Induction proves $r\equiv0$; no zero radius is used in the first-exit lemma.
For $T=0$, the conclusion is just $r_0=0$.

\par\noindent\textbf{Step 4: the failure event.}
The preceding argument proves the pathwise bound on
$\Omega_p\cap E_{T,\rm sp}(A_{\rm sp})\cap\mathcal S_{N,T,p}$.
Consequently any violation within the regularity event and the nonlinear stability event
belongs to $\Omega_p^c$, whose probability is at most $p$. No conditional
failure probability given the nonlinear stability event is asserted.
\end{proof}

\begin{corollary}[A simple sufficient horizon and the curvature-dependent rate]
\label{cor:fixed-schedule-high-prob-r-rate-v15}
Use the same model as $N$ varies, with fixed structural and anchor bounds and
uniform $K_0,B,A_{\rm sp}$. Fix $\upsilon\ge0$, $\zeta>0$ with $\upsilon+\zeta<1/4$, and $C_p,C_0,L_0>0$.
Suppose
\[
T\le N^{1/4-\upsilon-\zeta},\qquad p\ge N^{-C_p},\qquad
\NofJ\le L_0,\qquad 0\le\eta_\ell\le C_0,
\]
and choose $D$ as in \eqref{eq:def-supXY-v15}.
Define the controlled-curvature path condition $\mathcal R_N$ by
\begin{equation}
\label{eq:reference-response-event}
\mathcal R_N:=\left\{\int_0^T|\kappa_t^{Y,N}|\,dt<\infty,\quad
\mathfrak C_*^Y:=\sup_{t\le T}\mathfrak C_t^Y
\le C_0N^\upsilon(1+T),\quad
\mathfrak J_*^Y:=\sup_{t\le T}\mathfrak J_t^Y\le C_0N^\upsilon\right\}.
\end{equation}
For all sufficiently large $N$, the condition $\mathcal R_N$ implies the
nonlinear stability event $\mathcal S_{N,T,p}$ of \eqref{eq:smallness-v15}.
For every fixed $\varepsilon>0$, on
$\Omega_p\cap E_{T,\rm sp}(A_{\rm sp})\cap\mathcal R_N$,
\begin{equation}
\label{eq:r-rate-simplified-v15}
r_t\le C_\varepsilon\Xi^\varepsilon
\left[(N^{-1/2}+T/N)\mathfrak C_t^Y
+(N^{-1/4}+\sqrt{T/N})\mathfrak J_t^Y\right],\qquad t\le T.
\end{equation}
Thus the same bound holds with failure probability at most $p$ within
$E_{T,\rm sp}(A_{\rm sp})\cap\mathcal R_N$.
\end{corollary}
\begin{proof}
\noindent\textbf{Step 1: recall the deterministic envelopes.}
In this proof, write $\delta=\delta_{N,T,p}$ and
$\widetilde\delta=\widetilde\delta_{N,T,p}$ as defined in
\cref{eq:def-delta-v15,eq:def-tdelta-v15}.
Use \cref{lemma-uniform-sup-gamma-bound-convenient},
\eqref{eq:tilde-gamma-jump-time-final-rate}, and
\eqref{eq:Ccmp-localized-v15}. Uniform $K_0,B,A_{\rm sp}$ imply that
$K_{\rm pop},D$ are bounded by fixed powers of $\log\Xi$.
For \cref{lemma-uniform-sup-gamma-bound-convenient}, use total budget
$3p_0=3p/8<1$. Its population, radius, and concentration budgets are
then each $p_0$, and its scale agrees exactly with the present one:
\[
 \frac{3N(e+T)(\NofJ+1)(1+B)}{3p_0}
 =\frac{N(e+T)(\NofJ+1)(1+B)}{p_0}=\Xi.
\]
The polynomial-amplitude premise of that lemma holds with $c_M=1$.
Indeed, \cref{eq:population-exponential-amplitude} and the grid bound in
\cref{lem:population-moment-event} give
\[
 M_{\rm pop}\le\frac{4eQ}{p_0}
 \le\frac{C(e+T)(\NofJ+1)(1+B)}{p_0}
 =\frac{C\Xi}{N}\le C\Xi.
\]
For \eqref{eq:tilde-gamma-jump-time-final-rate} with budget $p_0=p/8$,
the present $\Xi\ge N(\NofJ+1)/p_0$ satisfies its scale condition
directly. For every $\theta>0$, the cited bounds give
\begin{equation}
\label{eq:main-rate-internal-envelopes}
\begin{aligned}
 \delta&\le C_\theta\Xi^\theta(N^{-1/2}+T/N),&
 \widetilde\delta&\le C_\theta\Xi^\theta N^{-1/2},\\
 c_D+c_J(D)q_p+\Lambda(D)&\le C_\theta\Xi^\theta.
\end{aligned}
\end{equation}
For $c_Jq_p$, apply the individual bounds with exponent $\theta/2$ to
each factor: $q_p=1+\sqrt{2\log(2(\NofJ+1)/p_0)}$ also satisfies
$q_p\le C_\theta\Xi^{\theta/2}$.

\par\noindent\textbf{Step 2: count the powers in the jump and curvature-response bounds.}
The amplitude bound $0\le\eta_\ell\le C_0$,
\eqref{eq:def-jump-amplification-v15}, and $\NofJ\le L_0$ imply
\[
 J_{N,T,p}\le(1+C_0c_J(D)q_p)^{L_0}
             \le C_\theta\Xi^{L_0\theta}.
\]
Using $\sqrt{a+b}\le\sqrt a+\sqrt b$ in
\eqref{eq:main-rate-internal-envelopes},
\[
 \sqrt\delta+\sqrt{\widetilde\delta}
 \le C_\theta\Xi^{\theta/2}(N^{-1/4}+\sqrt{T/N}).
\]
By \eqref{eq:def-B-v15} and the controlled-curvature path condition
\eqref{eq:reference-response-event},
\[
 \mathfrak B_*\le\delta\mathfrak C_*^Y
       +c_J(D)q_p(\sqrt\delta+\sqrt{\widetilde\delta})\mathfrak J_*^Y.
\]
Substitute the preceding estimates and $T/N\le N^{-1/2}$ to get
\[
 \mathfrak B_*
 \le C_\theta\Xi^\theta N^\upsilon(1+T)N^{-1/2}
       +C_\theta\Xi^{3\theta/2}N^\upsilon N^{-1/4}.
\]
Multiplying by $2J_{N,T,p}$ and using $\Xi\ge1$ yields
\[
 R\le C_\theta\Xi^{(L_0+3/2)\theta}N^\upsilon
         \{(1+T)N^{-1/2}+N^{-1/4}\}.
\]

\par\noindent\textbf{Step 3: verify the nonlinear stability event.}
Recall $\Xi=8N(e+T)(\NofJ+1)(1+B)/p$. The horizon and probability
assumptions imply $\Xi\le C N^{5/4+C_p}$. For the desired
$\varepsilon>0$, choose a fixed $\theta>0$ sufficiently small that
\begin{equation}
\label{eq:main-rate-internal-exponent}
 (5/4+C_p)(L_0+5/2)\theta\le\zeta/2,
 \qquad (L_0+3/2)\theta<\varepsilon.
\end{equation}
This choice depends only on the fixed parameters and $\varepsilon,\zeta$,
not on $N,T,p$. The preceding bounds and $N^\upsilon\ge1$ give
\[
\begin{aligned}
 T(\delta+c_DR)
 &\le C_\theta\Xi^{(L_0+5/2)\theta}N^\upsilon
       \left\{\frac{T(1+T)}{\sqrt N}+\frac{T}{N^{1/4}}\right\}\\
 &\le C_\theta N^{\zeta/2}
       (N^{-1/4-\zeta}+N^{-\upsilon-2\zeta}+N^{-\zeta})
 \le3C_\theta N^{-\zeta/2}\longrightarrow0.
\end{aligned}
\]
The second line substitutes $T\le N^{1/4-\upsilon-\zeta}$ into each term:
\[
 N^\upsilon T/\sqrt N\le N^{-1/4-\zeta},\qquad
 N^\upsilon T^2/\sqrt N\le N^{-\upsilon-2\zeta},\qquad
 N^\upsilon T/N^{1/4}\le N^{-\zeta}.
\]
The choice \eqref{eq:main-rate-internal-exponent} ensures
$(5/4+C_p)(L_0+5/2)\theta\le\zeta/2$.
For all sufficiently large $N$, the last bound is at most $1/4$, proving
$\mathcal R_N\subset\mathcal S_{N,T,p}$.

\par\noindent\textbf{Step 4: substitute into the tracking bound.}
The exact theorem and Steps 1--2 give
\[
 r_t\le C_\theta\Xi^{(L_0+3/2)\theta}
 \left[(N^{-1/2}+T/N)\mathfrak C_t^Y
 +(N^{-1/4}+\sqrt{T/N})\mathfrak J_t^Y\right].
\]
Since $(L_0+3/2)\theta<\varepsilon$, this proves
\eqref{eq:r-rate-simplified-v15}. The exact theorem continues to apply
at longer horizons on its nonlinear stability event \eqref{eq:smallness-v15}.
\end{proof}
